\section{Learning Outcomes}

\begin{frame}{Learning Outcomes}
    By the end of this session, you should be able to:
    \begin{itemize}\small\itemsep0.05em
        \item \textbf{Explain} what makes a small model strong in 2026, and what limits it on a device.
        \item \textbf{Compare} drafter-free speculation, pruning and low-rank methods by \emph{realised} speedup.
        \item \textbf{Describe} how diffusion LMs are built and sped up, and where they trail AR models.
        \item \textbf{Analyse} parametric memory: its capacity, and its freshness and attribution limits.
        \item \textbf{Explain} why stronger reasoners can hallucinate more, and choose an uncertainty estimator.
        \item \textbf{Evaluate} what probes, sparse autoencoders and steering vectors can and cannot show.
        \item \textbf{Assess} judges, reward models and leaderboards, which degrade when optimised.
        \item \textbf{Formulate} a research question from the October 2026 frontier.
    \end{itemize}
    \vspace{0.2em}
    {\scriptsize \textbf{Prerequisite knowledge:} speculative decoding; KV caching, GQA and sliding-window attention; prompt compression and KV-cache eviction; quantization, pruning and distillation, including on-policy distillation; mixture-of-experts, delta-rule linear attention and Titans; masked diffusion LMs (LLaDA); GRPO; reward models, LLM-as-a-judge and Chatbot Arena; binary grading and guessing, semantic entropy, Chain-of-Verification, fine-tuning on unknown facts; the refusal direction; chain-of-thought monitoring; agent time horizons; scaling laws and synthetic data.\par}
\end{frame}
